\chapter{What the bank would own and operate}
\label{ch:implementation}

The implementation has to answer a practical question: who can trust a decision,
and what can they inspect when the decision is challenged? A compliance
reviewer owns the meaning of a circular. A business owner sets the perimeter in
which that meaning will be used. Operations supplies and records the facts.
Engineering makes the approved meaning executable. A later incident may require
all four people to reconstruct one client or product decision.

This chapter follows that responsibility through one change. It starts with the
reviewer's view of the decision, then follows the evidence into a shared typed
model, tested RuleIR/Java and Catala targets, and a release. The detailed service
paths, storage rules and task history remain in Appendix~\ref{app:full-implementation-record}
of the technical companion; they are reference material for people who must
build or audit the route.

\section{Begin with the decision a person needs to make}

A reviewer opening the SPI example should see the source passage, the client's
financial and category evidence, the alternatives that were considered, and the
case that would change the answer. The reviewer should not have to reconstruct
that explanation from a list of fields. If the portfolio is missing but the
alternative net-asset route is established, the screen should say why the missing
value does not block this particular conclusion. If two records disagree, it
should show both records and identify the person who must resolve them.

The same view serves different people in different ways. Compliance needs the
source and adopted reading. Operations needs the action, owner and deadline.
Engineering needs the precise values and histories used to execute the rule. An
auditor needs to replay the result later. One maintained decision record can
support all four views if the reader starts with the conclusion and can open the
reason behind it.

\section{The shared frontend and the Catala routes}
\label{sec:implementation-shared-front-end}

Suppose the bank wants to calculate an amount from a client's holdings. The
earlier scalar rule can compare a supplied total with a threshold, but someone
must calculate that total first. A rule that selects eligible holdings and sums
their amounts needs the individual items as inputs. Moving the calculation
outside the rule would move its classification and completeness assumptions
outside the comparison too.

The shared frontend records the proposed calculation, its source passages,
scope, factual meanings and treatment of missing evidence in a
\emph{LegalRuleModel}. Each target translator receives that fixed model without
making another language-model call. It first checks which operations it
supports. The resulting \emph{capability report} either permits translation or
identifies the unsupported expression. A missing operation therefore remains
visible before a smaller program can be mistaken for the reviewed rule.

The two Catala routes serve different needs. The compatibility route translates
the supported RuleIR rules into Catala while keeping the existing decision,
explanation, evidence, transaction and release interfaces. The native route
translates the shared model directly. For the holdings calculation, a
\emph{record} groups each item's inclusion flag and amount; a \emph{list} holds
the records; filtering selects the included items; and summation adds their
amounts. The RuleIR target reports this collection model as unsupported, while
the native Catala target can execute it.

Two further types prevent information from being forced into an uninformative
flag. An \emph{option} represents either a present value or a known absence.
A \emph{payload variant} combines a category with the data appropriate to that
category, such as a tagged kind of holding with its associated amount. Neither
known absence nor a known category means that missing evidence has been found.
Version~1 already supports these richer values under a complete-input policy.
Version~2 adds scope conditions, explicit result states, defaults and exceptions,
exact scaling, separately chosen rounding, and typed calendar and collection
operations. It also supports pure reusable calculations, called helpers, and
several named results from one interpretation. Helper calls and type definitions
must have no circular dependencies.

Exactness and rounding answer different questions. Half of 101 monetary minor
units is 50.5 units. An exact scaling operation reports that the result is not
an integer number of units. A separately declared nearest rounding operation,
with ties away from zero, returns 51. Keeping those operations distinct prevents
the compiler from choosing a monetary policy implicitly. Calendar operations
likewise have named conventions: adding a month with month-end clamping uses
the final day of the destination month when the original day does not exist.

The input policy is part of the meaning. The \texttt{complete.v1} profile
requires every declared input to be usable. Version~2 also supports
\texttt{partial.v1}, which follows the observations actually used in the
calculation. Consider the retained synthetic example of a record whose amount
is known to be one-third but whose inclusion flag is unknown. Reading the
amount returns one-third. Returning the whole record gives a partial result
with the two field states preserved. Filtering on the inclusion flag leaves
membership unresolved, so the sum is unknown. A complete empty list, by
contrast, sums to zero; an incomplete list cannot establish that no further
items exist.

RuleIR checks for conflicts across the facts in a rule's static dependency set
before evaluating its expression. The partial Catala policy can inspect an
unaffected field without inspecting another field's conflict. These are different
policies, so the RuleIR target rejects \texttt{partial.v1}. The compatibility
route retains RuleIR's original policy. In both native profiles, Python
validates and encodes observations and decodes results; Catala performs the
arithmetic, selection and propagation of known, unknown, conflict and other
declared result states.

\begin{table}[htbp]
\centering
\caption{The executable layers used by the completed comparison.}
\label{tab:shared-executable-layers}
\begin{tabularx}{\textwidth}{P{34mm}YY}
\toprule
Layer & Responsibility & Evidence boundary\\
\midrule
Shared frontend & Retain sources and scope; record one typed interpretation, evidence policy and outputs & Common model and interpretation hashes\\
\addlinespace[0.4em]
RuleIR target & Execute the bounded scalar and finite-duty profile through Python and Java & Existing RuleIR semantics, traces and release checks\\
\addlinespace[0.4em]
Catala target & Lower the common profile or emit richer native Catala with declared policies & Catala build, interpreter, Java and exact-reference checks\\
\addlinespace[0.4em]
Host boundary & Validate facts and evidence, decode outputs and preserve provenance & Input, result, package and transaction invariants\\
\bottomrule
\end{tabularx}
\end{table}

Freezing the model makes a disagreement introduced by translation easier to
isolate. A missing source qualification can still be present in that common
model. The source checks in Chapter~\ref{ch:ensemble} address this earlier
opportunity for error.

\subsection{What the completed checks establish}
\label{sec:implementation-catala-evidence}

The retained evidence has several scopes. The modular replay sent the same
frozen model and evidence policy through both targets for 44 paired cases. Each
pair preserved the model and interpretation hashes, and each compiled result
matched the original RuleIR evaluator's complete semantic trace. Seven richer
outputs executed as native Catala, uninstrumented Java and exact checks inside
the pinned Catala interpreter; the RuleIR target reported those models as
unsupported. These checks establish finite translation and execution behavior
for the declared controls. They do not establish that the source readings are
legally correct.

Version~2 received a separate set of 91 exact checks for the extended semantics.
Section~\ref{sec:evaluation-catala-comparison} distinguishes these checks from
the earlier backend tests and source-generation study.
Appendix~\ref{app:catala-reproduction} in the technical companion identifies
the commands, source versions, logs and detailed regression accounting.

The common fragment retained its checked behavior, and the native target
executes the documented richer models. Whether that additional expressiveness
improves interpretation or review is a separate empirical question. The
completed engineering tests do not supply a performance or legal-quality
ranking, human review evidence or production approval.

\subsection{The bounds that remain part of the interface}
\label{sec:implementation-catala-bounds}

The native profile is a closed typed vocabulary with acyclic helper and type
graphs. The generated interface enforces declared limits on native types,
expression depth, inputs and outputs, candidate source size, observations and
expanded types. Imported modules outside the pinned library set, recursive
types, arbitrary Catala operators and unbounded collections remain unsupported.
All-output execution now evaluates the sealed program once and returns each
output with the same value, status and evidence as an individual call. A resource
report exposes the type and depth expansion introduced by state encoding before
compilation. The earlier directly
authored native programs could raise an execution failure for conflicting
consequences. That historical behavior must be distinguished from the completed
shared translator, which encodes strict exceptions and structured result states.
Two active exceptions conflict there even when their consequences agree.

The current native trace records branches emitted into Java and their source
positions. Observation hooks are added after the ordinary compiler passes, so
both the original and instrumented compilers run the invariant checks, including
the nested-scope check that the first tracing implementation could not pass.
Standard-library internals and source choices removed by optimization are not
observed. Comparing plain Java, instrumented Java and the original Catala
interpreter helps expose changes introduced by tracing, but those finite checks
do not prove that the instrumentation preserves every program.

These limits are useful precisely because they are visible. An unsupported
construct produces a capability report or a reviewable build failure. It is not
silently rewritten as a smaller rule. A bank release would still require the
source, interpretation, evidence mappings, identity controls and approvals
described later in this chapter.

\section{Follow one change from source to release}

The route is easier to understand as a sequence of questions than as a software
architecture. What did the circular say? Which activity and legal entity are in
scope? Which reading did the bank adopt? What facts does that reading require?
Which cases would expose an error? Does the program return the same result as the
reviewed rule? Who approved this exact version, and when may the bank use it?

\begin{figure}[htbp]
\centering
\begin{tikzpicture}[node distance=7mm and 7mm,
 box/.style={draw=teal,fill=pale,rounded corners,align=center,text width=35mm,minimum height=15mm,font=\small},
 arr/.style={-{Latex[length=2mm]},thick,draw=navy}]
\node[box] (source) {Source passage and\\ business scope};
\node[box,right=of source] (meaning) {Reviewed meaning\\ and open questions};
\node[box,right=of meaning] (facts) {Facts, cases and\\ evidence owners};
\node[box,below=of facts] (program) {Reference rule and\\ generated Java};
\node[box,left=of program] (approval) {Independent checks\\ and approvals};
\node[box,left=of approval] (use) {Supervised use\\ and replay};
\draw[arr] (source)--(meaning); \draw[arr] (meaning)--(facts);\draw[arr] (facts)--(program);
\draw[arr] (program)--(approval); \draw[arr] (approval)--(use);
\draw[arr] (approval.north)--(meaning.south);
\end{tikzpicture}
\caption{A bank release carries a reviewed meaning and its evidence through
execution. Each step answers a different question.}
\label{fig:reader-implementation-route}
\end{figure}

The prototype now exercises this route for public documents and synthetic facts.
Consider the gift review. The reviewer begins with the selected paragraph and
source edition, sees alternative readings and missing coverage, and examines a
concrete distinguishing case where the fact meanings can be aligned. Separate
control components keep their answers and open questions when Java is generated.
If the source changes, affected work is reconsidered while the earlier decision
remains reproducible.

That journey includes different kinds of progress. Retrieving the FAQ supplies
context. Running a case exposes a program difference. An authorized decision
adopts a meaning for the bank. The first two are available in local development;
the third requires the institution's own authority and evidence. The demonstrated
route therefore supplies draft controls for review, with the open questions
carried through to delivery.

\section{Keep the release tied to its meaning}

A tested program is only useful when the bank can identify the rule it implements.
The release therefore carries the source edition, the adopted interpretation,
the factual definitions, the supported dates and activities, the test evidence
and the exact Java package together. Changing one of those parts creates a new
version that needs its own review. A content identifier helps detect accidental
mixing of versions; it does not make an administrator, source or interpretation
authoritative.

The same principle governs a correction. If a client withdraws consent on
Tuesday, a Monday decision must remain reproducible with the information known
on Monday. A Wednesday amendment must show which cases and clients change. The
bank can replace a defective software package while preserving the legal version
that applied to an earlier transaction. Historical replay is therefore part of
the control's meaning, not a convenience added after deployment.

\section{What the prototype has and has not demonstrated}

The engineering record contains checks for selected rules, edge cases,
deliberate mutations, release identity, withdrawal, amendment and archive
restoration. The interpretation route now also generates alternatives and
checks their source support. Earlier rounds included live development calls;
later integration exercises used those retained responses and controlled
inputs, with unresolved interpretations retained through delivery.

The distinction matters for the next investment decision. The mechanism can now
be exercised, but no comparative study has measured how often it misses a
requirement on unfamiliar material, how quickly a reviewer can use its reports,
or how it behaves within the bank's identity and data controls. Public sources
and synthetic identities remain the basis of the local demonstrations.

Those unanswered questions belong to the evaluation and operating decisions in
the last two chapters. Keeping them separate prevents a successful build from
being mistaken for a legal approval or a production control.

\section{The boundary before a bank pilot}

Before a supervised pilot, the bank would have to name the legal entities,
products, client categories, source families and decision points included in the
study. It would need an authorised legal owner, data owners for each fact,
independent engineering and compliance review, and a response for unknown or
conflicting evidence. The pilot should run beside the existing process so that
staff can compare decisions and record material disagreements. A case that cannot
be supported should stop or go to a named reviewer; it should not be converted
into a plausible default.

The implementation record makes the route concrete. The management decision is
whether the bank is willing to supply the authority, evidence and supervision
that the current public demonstration cannot supply by itself.
